\subsection{CHARACTERIZATION OF PROPER MAPPINGS BY COMPACTNESS PROPERTIES}

In this subsection we shall denote by P a space consisting of a single point, with its unique topology.

Lemma 1. Let X be a topological space such that the constant mapping \(X \rightarrow P\) is proper. Then X is quasi-compact.

(We shall see a little later on (Theorem 1, Corollary 1) that this property characterizes quasi-compact spaces.)

We may restrict ourselves to the case where X is not empty. Let F be a filter on X, and \(X' = X \cup \{\omega\}\) the topological space associated with F (§ 6, no. 5, Example). Let \(\Delta\) be the subset of \(X \times X'\) consisting of all \((x, x)\) where \(x \in X\) , and let \(F = \overline{\Delta}\) be the closure of \(\Delta\) in \(X \times X'\) . In view of the hypothesis on X, the image of F under the projection \(X \times X' \to X'\) is closed in \(X'\) ; this image contains X and therefore contains \(\omega\) , which lies in the closure of X; in other words, there is a point \(x \in X\) such that \((x, \omega) \in F\) . By the definition of the topology of \(X \times X'\) , this means that, for each neighbourhood V of x in X and each \(M \in F\) , we have \((V \times M) \cap \Delta \neq \emptyset\) , i.e.~\(V \cap M \neq \emptyset\) , so that x is a cluster point of the filter F, and therefore X is quasi-compact.

Q.E.D.

THEOREM I. Let \(f: \mathbf{X} \to \mathbf{Y}\) be a continuous mapping. Then the following four statements are equivalent:

\begin{enumerate}
\def\labelenumi{\alph{enumi})}
\item
  \(f\) is proper.
\item
  \(f\) is closed and \(\overline{f}^1(y)\) is quasi-compact for each \(y \in \mathbf{Y}\) .
\item
  If \(\mathfrak{F}\) is a filter on \(\mathbf{X}\) and if \(y\in \mathbf{Y}\) is a cluster point of \(f(\mathfrak{F})\) then there is a cluster point \(x\) of \(\mathfrak{F}\) such that \(f(x) = y\) .
\item
  If \(\mathfrak{U}\) is an ultrafilter on \(\mathbf{X}\) and if \(y\in \mathbf{Y}\) is a limit point of the ultrafilter base \(f(\mathfrak{U})\) , then there is a limit point \(x\) of \(\mathfrak{U}\) such that \(f(x) = y\) .
\item
  \(\Longrightarrow\) b): If \(f\) is proper then \(f\) is closed (no. 1, Proposition 1) and for each \(y \in Y\) the mapping \(f_{\{y\}}: \overline{f}^1(y) \to \{y\}\) is proper {[}no. 1, Proposition 3a){]}. By Lemma 1, this implies that \(\overline{f}^1(y)\) is quasi-compact.
\item
  \(\Longrightarrow\) c): Suppose \(\mathfrak{F}\) and \(y\) satisfy the hypotheses of c). Let \(\mathfrak{B}\) be the filter base on X formed by the closures of the sets of \(\mathfrak{F}\) . Since \(f\) is closed, we have \(f(\overline{\mathbf{M}}) = \overline{f}(\mathbf{M})\) for each \(\mathbf{M} \in \mathfrak{F}\) (§5, no. 4, Proposition 9). This shows that the sets \(\overline{\mathbf{M}} \cap \overline{f}^1(y)\) are non-empty for all \(\mathbf{M} \in \mathfrak{F}\) , and hence form a filter base on \(\overline{f}^1(y)\) whose elements are closed subsets of \(\overline{f}^1(y)\) . Since \(\overline{f}^1(y)\) is quasi-compact, there is a point \(x \in \overline{f}^1(y)\) which belongs to all the sets \(\overline{\mathbf{M}}\) as M runs through \(\mathfrak{F}\) . Hence \(f(x) = y\) and \(x\) is a cluster point of \(\mathfrak{F}\) .
\end{enumerate}

\begin{enumerate}
\def\labelenumi{\alph{enumi})}
\setcounter{enumi}{2}
\tightlist
\item
  \(\Longrightarrow\) d): Trivial.
\end{enumerate}

\begin{enumerate}
\def\labelenumi{\alph{enumi})}
\setcounter{enumi}{3}
\tightlist
\item
  \(\Longrightarrow\) a): We show first that if d) is satisfied, then \(f\) is a closed mapping. Let A be a non-empty closed subset of X and let \(\mathfrak{F}\) be the filter of subsets of X which contain A. Then A is the set of cluster points of \(\mathfrak{F}\) . Let B be the set of cluster points of the filter base \(f(\mathfrak{F})\) on Y; B is closed and clearly contains \(f(A)\) ; we shall show that \(B = f(A)\) . Let \(y \in B\) and let \(\mathfrak{B}\) be the neighbourhood filter of \(y\) in Y; then by hypothesis every set of \(\mathfrak{W} = \overline{f}^{1}(\mathfrak{V})\) meets every set of \(\mathfrak{F}\) ; hence \(\mathfrak{W}\) is a filter base on X and there is an ultrafilter \(\mathfrak{U}\) on X which is finer than both \(\mathfrak{F}\) and the filter whose base is \(\mathfrak{W}\) (§ 6, no. 2, Proposition 1, Corollary 1 and no. 4, Theorem 1). The ultrafilter whose base is \(f(\mathfrak{U})\) is finer than \(\mathfrak{V}\) and therefore converges to \(y\) . By virtue of d) there is a point \(x \in X\) such that \(f(x) = y\) and \(\mathfrak{U}\) converges to \(x\) ; since \(\mathfrak{U}\) is finer than \(\mathfrak{F}\) , \(x\) is a cluster point of \(\mathfrak{F}\) ; hence \(x \in A\) . This shows that \(B = f(A)\) and therefore that \(f\) is closed.
\end{enumerate}

To complete the proof we have to show that \(f \times \iota_{\mathbf{Z}}\) is closed for every topological space \(\mathbf{Z}\) . From what has been proved it is enough to show that if \(f\) satisfies condition d), then so does \(f \times \iota_{\mathbf{Z}}\) . This is a consequence of the following general lemma:

Lemma 2. If \((f_{\iota})_{\iota\in I}\) is a family of continuous mappings \(f_{\iota}:\mathbf{X}_{\iota}\to \mathbf{Y}_{\iota}\) each of which satisfies condition d) of Theorem 1, then the product mapping

\[
f: (x _ {\iota}) \rightarrow (f _ {\iota} (x _ {\iota}))
\]

also satisfies d).

Let \(\mathfrak{U}\) be an ultrafilter on \(\mathbf{X} = \prod_{\iota} \mathbf{X}_{\iota}\) , and let \(y = (y_{\iota})\) be a point of \(\mathbf{Y} = \prod_{\iota} \mathbf{Y}_{\iota}\) such that \(f(\mathfrak{U})\) converges to \(y\) . This means that each of the ultrafilter bases \(\operatorname{pr}_{\iota}(f(\mathfrak{U})) = f_{\iota} (\operatorname{pr}_{\iota}(\mathfrak{U}))\) converges to \(y_{\iota}\) ( \(\S 7\) , no. 6, Proposition 10, Corollary 1). By virtue of condition d), for each \(\iota \in I\) there exists \(x_{\iota} \in \mathbf{X}_{\iota}\) such that \(f_{\iota}(x_{\iota}) = y_{\iota}\) and \(\operatorname{pr}_{\iota}(\mathfrak{U})\) converges to \(x_{\iota}\) ; but then \(\mathfrak{U}\) converges to \(x = (x_{\iota})\) (loc. cit.) and we have \(f(x) = y\) . This completes the proof of Lemma 2 and hence of Theorem 1.

COROLLARY 1. A topological space X is quasi-compact if and only if the mapping \(X \rightarrow P\) is proper.

Apply a) \(\Longleftrightarrow\) b) to \(\mathbf{X} \to \mathbf{P}\) .

COROLLARY 2. Every continuous mapping f of a quasi-compact space X into a Hausdorff space Y is proper.

The composition \(X \xrightarrow{f} Y \to P\) is proper by Corollary 1; hence f is proper by no. 1, Proposition 5 d). Alternatively we may apply the criterion b) of Theorem 1, using § 9, no. 4, Theorem 2, Corollary 2.

COROLLARY 3. If \((f_{\iota})\) is a family of proper mappings, then the product mapping \((x_{\iota})\to (f_{\iota}(x_{\iota}))\) is proper.

In view of Theorem 1, this is just Lemma 2 above.

If we apply this corollary to the family of mappings \(X_{\iota} \rightarrow P\) and use Corollary 1, we have Tychonoff's theorem (§ 9, no 5, Theorem 3).

COROLLARY 4. Let X be a Hausdorff space, and let \(f_{\iota}: \mathbf{X} \to \mathbf{Y}_{\iota}\) be a family of proper mappings. Then the mapping \(f: x \to (f_{\iota}(x))\) of X into \(\prod_{\iota} Y_{\iota}\) is proper.

The proof is the same as in the case of a finite family (no. 1, Proposition 5, Corollary 3), using Corollary 3 above and the fact that the diagonal of \(\mathbf{X}^{\mathbf{I}}\) is closed ( \(\S\) 8, no. 1, Proposition 1).

COROLLARY 5. If X is any quasi-compact space and Y is any topological space, then the projection \(\mathbf{pr}_2: \mathbf{X} \times \mathbf{Y} \to \mathbf{Y}\) is proper.

For we may identify Y with \(P \times Y\) and \(pr_{2}\) with the product of \(X \rightarrow P\) and \(\iota_{Y}\) , both of which are proper mappings.

Example. Let X be a set, and let \(f: \mathbf{X} \to \mathbf{X}'\) be a mapping of X onto a topological space \(\mathbf{X}'\) ; topologize X with the inverse image under \(f\) of the topology of \(\mathbf{X}'\) . Then \(f\) is proper, for \(f\) is closed ( \(\S 5\) , no. 1, Example 3) and the inverse image of a point of \(\mathbf{X}'\) is a subspace of X whose topology is the coarsest topology and is therefore quasi-compact.

Remark. When Y is Hausdorff, condition d) of Theorem 1 is equivalent to the following:

d') If \(\mathfrak{U}\) is an ultrafilter on \(\mathbf{X}\) such that \(f(\mathfrak{U})\) is a convergent filter base, then \(\mathfrak{U}\) is convergent.

For if \(\mathfrak{U}\) converges to \(x\) and \(f(\mathfrak{U})\) converges to \(y\) , then the uniqueness of the limit in \(\mathbf{Y}\) and the continuity of \(f\) show that we must have \(y = f(x)\) . Likewise, Y being Hausdorff, condition c) of Theorem 1 is equivalent to:

c') If \(\mathfrak{F}\) is a filter on \(\mathbf{X}\) such that \(f(\mathfrak{F})\) has a cluster point, then \(\mathfrak{F}\) has a cluster point.

\[
\text { For } \quad \mathrm{c}) \Longrightarrow \mathrm{c} ^ {\prime}) \Longrightarrow \mathrm{d} ^ {\prime}) \Longrightarrow \mathrm{d}) \Longrightarrow \mathrm{c}).
\]

On the other hand, if Y is not Hausdorff, then d') no longer implies d); for example, take X to consist of one point and Y to consist of two points, with the coarsest topology.

PROPOSITION 6. Let \(f: \mathbf{X} \to \mathbf{Y}\) be a proper mapping, and let \(\mathbf{K}\) be a quasi-compact subset of \(\mathbf{Y}\) . Then \(f^1(\mathbf{K})\) is quasi-compact.

By Proposition 3 of no. 1, the mapping \(f_{\mathbf{K}}: \overline{f}^{1}(\mathbf{K}) \to \mathbf{K}\) is proper. Since \(\mathbf{K} \to \mathbf{P}\) is a proper mapping (Theorem 1, Corollary 1) it follows from no. 1, Proposition 5 a) that the composition \(\overline{f}^{1}(\mathbf{K}) \xrightarrow{f_{\mathbf{K}}} \mathbf{K} \to \mathbf{P}\) is proper, whence \(\overline{f}^{1}(\mathbf{K})\) is quasi-compact by Theorem 1, Corollary 1.

